% Exercice: Exercice 11
% Sujet: Algèbre
% Généré automatiquement

\begin{exercice}[Algèbre]
\label{ex:3}

1^\{1\} points\\
PROBLEME\\
Le problème comporte deux parties Aet B.\\
\item A) La courbe
r), ci-contre est la représentation graphique dune fonction numérique\\
dans un repère orthonormé $(0,5,1)$\\
Par lecture graphique\\
Déterminer l'ensemble de définition Df de\\
1^\{7\} et en 1*\\
,5 pt\\
ainsules $limites en$\\
2- Préciser le sens de variation de\\
dans\\
Résoudre\\
les inéquations\\
\item i) f(x)< 0
f(x)>0.\\
1pt\\
Déterminer f(-1) ; f(0) ; f(-1) où f est la\\
0,7^\{5\} pt\\
fonction derivée de\\
0,7^\{5\} pt\\
5- Dresser\\
tableau de variation de f\\
1l) On suppose que pour tout réel\\
4^\{1\}\\
$f(x)=ax +b+$\\
où a, b et c sont 3 réels.\\
En utilisant la question\\
montrer que\\
réels a, b et\\
vérifient le système\\
2b + c\\
0,7^\{5\} pt\\
=-2\\
2- Choisir la lettre correspondant\\
la bonne\\
0,7^\{5\} pt\\
reponse\\
Le triplet (a\\
\item c) est egal a
\item e) $(-2,1,8)$ ; 1) $(-2,4,-2)$ ; 9) $(1,2,4)$ ; h) $(0,-1, 0)$

\end{exercice}

% --- Fin Exercice 11 ---
